\documentclass[../veristore_security_model.tex]{subfiles}

\begin{document}
\section{Data Flow Diagram}

% COMMENT: This section contains the graphical representation of data flows through veri-store.
% The DFD shows all major components, entry points, data stores, and how information moves
% between them. Use different colors for trust levels: green for trusted processes, red for
% untrusted entities/data, orange for semi-trusted. Number all data flows sequentially and label
% them with concise descriptions.

\begin{figure}[ht!]
\centering
\begin{tikzpicture}[node distance=2.5cm]

% External entities
\node (client) [entity, fill=trusted!30] {Client};
\node (network) [entity, fill=untrusted!30, right of=client, xshift=3cm] {Network};
\node (storage) [datastore, fill=semitrusted!30, below of=network, yshift=-1cm] {Fragment \newline Storage};

% Processes (all trusted - our code)
\node (encode) [process, below of=client, yshift=-0.5cm] {Encode \newline Block};
\node (fingerprint) [process, right of=encode, xshift=2cm] {Compute \newline Fingerprint};
\node (distribute) [process, below of=fingerprint, yshift=-0.5cm] {Distribute \newline Fragments};
\node (verify) [process, below of=distribute, yshift=-0.5cm] {Verify \newline Fragments};
\node (reconstruct) [process, left of=verify, xshift=-2cm] {Reconstruct \newline Block};

% Data flows
\draw [arrow] (client) -- node[anchor=east] {1. Data block} (encode);
\draw [arrow] (encode) -- node[anchor=south] {2. Fragments} (fingerprint);
\draw [arrow] (fingerprint) -- node[anchor=west] {3. Fingerprints} (distribute);
\draw [arrow] (distribute) -- node[anchor=west] {4. Fragment + \newline fingerprint} (network);
\draw [arrow] (network) -- node[anchor=west] {5. Store} (storage);
\draw [arrow] (storage) -- node[anchor=east] {6. Retrieve} (verify);
\draw [arrow] (verify) -- node[anchor=south] {7. Verified \newline fragments} (reconstruct);
\draw [arrow] (reconstruct) -- node[anchor=east] {8. Reconstructed \newline block} (client);

% Cross-checksum verification (dotted line to show verification flow)
\draw [arrow, dashed] (storage) to[bend left=30] node[anchor=west] {9. Cross-check \newline fingerprints} (verify);

\end{tikzpicture}
\caption{High-level data flow diagram for veri-store store and retrieve operations.}
\label{fig:dfd}
\end{figure}

\subsection{DFD Component Descriptions}

% COMMENT: Provide brief descriptions of each component in the DFD. Explain what each process
% does, what trust level it operates at, and what security properties it maintains. Describe
% external entities and data stores. Note any assumptions about data flows (e.g., authenticated
% channels, encrypted storage).

\subsubsection{Client}

% COMMENT: External entity representing the user or application storing/retrieving data. Operates
% at "authenticated client" trust level. Assumed to be benign but may have bugs.

\subsubsection{Encode Block}

% COMMENT: Trusted process that implements Reed-Solomon erasure coding. Takes original block,
% produces n fragments such that any m can reconstruct the original. Must correctly implement
% the encoding algorithm to ensure fragments are valid.

\subsubsection{Compute Fingerprint}

% COMMENT: Trusted process that generates homomorphic fingerprints for each fragment using division
% fingerprinting over F_256. Must use cryptographically secure random parameter r and correctly
% compute fingerprints to ensure collision resistance.

\subsubsection{Distribute Fragments}

% COMMENT: Trusted process that sends fragments and fingerprints to storage servers. Handles
% network communication and retries. May need to authenticate servers to prevent MITM attacks.

\subsubsection{Network}

% COMMENT: Untrusted external entity representing network infrastructure. Adversary may observe,
% modify, or inject messages. Protection via authentication (TLS) or message authentication codes
% may be needed.

\subsubsection{Fragment Storage}

% COMMENT: Semi-trusted data store representing persistent storage on servers. Stores fragments
% and fingerprints. Servers may fail or return corrupted data but are detected via verification.
% Storage is assumed to be persistent but not necessarily confidential.

\subsubsection{Verify Fragments}

% COMMENT: Trusted process that validates fragment integrity using homomorphic fingerprints.
% Retrieves fragments and fingerprints from storage, performs cross-checksum verification to
% detect corrupted or forged fragments. Rejects invalid fragments.

\subsubsection{Reconstruct Block}

% COMMENT: Trusted process that decodes original block from m valid fragments using Reed-Solomon
% decoding. Only operates on verified fragments to ensure integrity of reconstructed data.
\end{document}